% ----------------------------------------------------------
% Introdução
% ----------------------------------------------------------
\chapter{Introdução}

O sistema prisional brasileiro vem sofrendo diariamente com problemas recorrentes e desafiadores, que afetam diretamente de forma estrutural e operacional o convívio de cada envolvido, indo contra a finalidade ressocializadora prevista na Lei de Execução Penal (LEP) . Recentemente, segundo dados do Sistema de Informações do Departamento Penitenciário Nacional(SISDEPEN), o Brasil possui uma população carcerária com cerca de  909.067 pessoas, referente ao segundo semestre de 2024, sendo 674.016 em celas físicas, o que representa um aumento de quase 6\% em dois anos. Assim, diante deste cenário, faz-se necessário maior rigor na gestão das penas e na aplicação dos benefícios legais, evitando distorções que possam impactar tanto a segurança pública quanto os direitos individuais de cada detento.\cite {senappen2024, lei12714}

Entre os fatores envolvidos que favorecem este cenário e comprometem a efetivação da execução penal, destaca-se a falha no cálculo manual da pena, que é realizado durante o processo de execução criminal.  Essa operação por fazer uso de cálculos complexos que necessitam da interpretação da Lei de execução penal(LEP), torna-se sujeita a falhas que passam despercebidas durante o seu procedimento. Um exemplo é o caso ocorrido na Defensoria Pública do Ceará, em que um réu permaneceu preso mesmo após cumprir integralmente a pena, evidenciando as consequências que erros nesse processo podem causar ao acusado.\cite{lep, defensoriaCE}

Com efeito, um cálculo incorreto pode acarretar em uma prisão prolongada além do previsto, resultando em uma permanência do apenado no sistema prisional, ou caso contrário, uma soltura prematura indevida, fora do tempo determinado segundo a lei. Consequentemente, esses erros afetam a gestão de vagas nas penitenciárias e a finalidade ressocializadora da pena, gerando a necessidade de maior precisão e confiabilidade no processo. \cite{defensoriaCE}

Um cálculo correto e preciso não garante apenas que o apenado cumpra o tempo adequado de execução penal, mas também assegura a legalidade do cumprimento da pena e a efetivação da ressocialização prevista na Lei de Execução Penal. Garantir que o processo seja realizado corretamente, desde o trânsito em julgado até a extinção da pena, é essencial para evitar distorções que comprometam direitos individuais e a segurança jurídica. Nesse sentido, torna-se necessário adotar mecanismos que reduzam falhas humanas e promovam maior eficiência na gestão das penas.\cite{lep}

A calculadora de execução criminal é de suma importância, pois busca garantir transparência, eficácia e precisão quanto ao cumprimento das penas, favorecendo o apenado, o sistema de justiça e a sociedade como um todo.  Possibilitando ao apenado e sua defesa, uma compreensão clara e objetiva de sua situação penal, fazendo-os enxergar onde precisam de maior atenção, quanto às normas a serem cumpridas, quando podem solicitar progressão de regime, o que resta para chegar na fase liberdade e ficando seguros acerca do trâmite necessário. Promove a prevenção quanto a falhas que podem afetar diretamente o cálculo de uma execução penal. Por fim, visa garantir que, compreendendo e executando o processo de forma coesa e justa, de acordo com o que as normas regem e os resultados dos cálculos transmitem o necessário a se cumprir, reflete na retomada e reinserção do apenado à sociedade.

\section{Objetivo Geral}

Este trabalho tem como objetivo garantir que a automação do cálculo de execução penal seja realizada com precisão e eficiência, conforme os parâmetros estabelecidos pela LEP.

\subsection{Objetivos Específicos}
\begin{itemize}
\item Facilitar o processo de cálculo de execução criminal acerca das regulamentaçõe previstas pela lei;
\item Garantir que o sistema automatizado preserve os benefícios previstos na Lei de Execução Penal, como progressão de regime e remição por estudo ou trabalho, contribuindo para a eficiência do sistema carcerário e para a finalidade ressocializadora da pena;
\item Implementar funcionalidades de progressão, remição e detração de pena;
\item Avaliar os impactos da automatização na redução de falhas e na otimização dos trâmites burocráticos.
\end{itemize}
